% EM-FLAT upload variant for Elsevier Editorial Manager (no subfolders): sections inlined,
% bibliography inlined from .bbl, figure paths bare. Upload this .tex as Manuscript and the
% three PNGs as Figure items at the same folder level.
% IPIS Module 4 -- composed coverage certificate. Computers & Chemical Engineering (Elsevier).
% Build: pdflatex main && bibtex main && pdflatex main && pdflatex main
\documentclass[review,12pt]{elsarticle}

\usepackage{amsmath,amssymb}
% elsarticle predefines \proof; free it so amsthm can provide the proof environment
\let\proof\relax
\let\endproof\relax
\usepackage{amsthm}
\usepackage{booktabs}
\usepackage{array}
\usepackage{graphicx}
\usepackage{lineno}
\usepackage{url}

% theorem environments (structure unchanged from the reviewed version)
\newtheorem{theorem}{Theorem}
\newtheorem{corollary}{Corollary}
\newtheorem{lemma}{Lemma}
\theoremstyle{definition}
\newtheorem{assumption}{Assumption}
\newtheorem{definition}{Definition}

\journal{Computers \& Chemical Engineering}

\begin{document}
\begin{frontmatter}

\title{A composed coverage certificate for closed-loop process operation:
certified joint product-quality and equipment-survival guarantees under feedback}

\author{Bien Busico}
\address{Chemical Engineer, Quezon City, Philippines}
\ead{bienbusico@gmail.com}

\begin{abstract}
A self-optimizing plant is safe only when a joint event holds: the product stays in
specification while the equipment survives to its next maintenance window. We show that this
joint guarantee can be certified in closed loop, that feedback does not void it, and that a
single convex constraint forces a profit-seeking optimizer to respect it. Soft sensors,
prognostic models, and real-time optimization (RTO) each now ship distribution-free coverage
guarantees, yet the joint event that determines safety follows from none of them in
isolation, and feedback is widely believed to break the exchangeability on which conformal
validity rests. We derive a per-cycle \emph{composed} coverage certificate for the joint
event, with terms that map one-to-one onto the three modules: the soft sensor's miscoverage,
the prognostic bound's miscoverage, a similitude-departure penalty contributed by each
module's regime shift, and a selection penalty. A causal-timing argument then shows that the
selection penalty is exactly zero whenever the decision is formed before the measurement it
is conditioned against, which dissolves the feedback objection. A corollary collapses the
certificate into a single convex similitude-departure budget that a modifier-adaptation RTO
enforces directly, so the optimizer cannot push economics hard enough to invalidate its own
sensor and life guarantees, and the flagship guarantee and the operating floor become the
same constraint. A negative control isolates the mechanism. On a calibrated debutanizer
digital twin running the real soft-sensor, prognostic, and RTO implementations, the certified
floor $1-(\alpha_1+\alpha_2)-\varepsilon = 0.75$ is met empirically, with joint safety
coverage $0.988$ ($95\%$ interval $[0.976, 0.994]$) when the budget is enforced; removing the
budget, the single causal ingredient, collapses joint coverage to $0.000$. Reflux-pump life
carries the entire contrast, with remaining-life coverage of $1.000$ against $0.013$. None of
the three guarantees implies system safety on its own; the composition, and the convex budget
that enforces it, is what the certificate provides.

\end{abstract}

\begin{keyword}
conformal prediction \sep real-time optimization \sep soft sensor \sep
predictive maintenance \sep digital twin \sep distillation \sep
uncertainty quantification \sep process safety
\end{keyword}

\end{frontmatter}
\linenumbers

\section{Introduction}
\label{sec:intro}

Modern process operation increasingly layers data-driven components on top of the base
control system. Soft sensors infer product qualities that are too slow or too costly to
measure online; prognostic models estimate the remaining useful life of critical
equipment; and real-time optimization (RTO) continually re-tunes setpoints toward the
economic optimum. Each of these layers has, on its own, acquired rigorous uncertainty
quantification. Conformal prediction in particular supplies distribution-free,
finite-sample coverage guarantees~\citep{vovk2005,shafer2008,angelopoulos2023}, and
calibrated intervals are now practical for soft sensing, for prognostics, and for
constraint handling in optimization.

Safe operation, however, is not a property of any one layer. It is a property of the
joint outcome: the product must be in specification and the equipment must survive to its
next maintenance window. The optimizer acts on estimates, a soft-sensor interval and a
remaining-life bound, yet the event that determines safety is defined on the truths it
cannot observe. The coupling is tighter still, because the optimizer's own decision moves
every layer's operating regime at once. The reflux lever it raises for economic gain
simultaneously shifts the soft sensor away from its calibration regime and loads the pump
whose life the prognostic model is predicting. The three guarantees are therefore coupled
through the decision, and the guarantee a closed-loop operator actually needs, coverage of
the joint safety event under this coupling, does not follow from the three guarantees
taken in isolation.

A second difficulty compounds the first: conformal validity is widely held to be
incompatible with feedback. Conformal coverage rests on exchangeability between the
calibration data and the test point, and a closed loop breaks it, because the test point's
operating regime is selected by the controller using the very predictions under test. The
usual conclusion is that conformal guarantees do not survive in the loop, so conformal
methods are either confined to open-loop monitoring or deployed in feedback without a
guarantee.

This paper resolves both difficulties for the integrated system. Its contributions are:
\begin{enumerate}
\item A per-cycle \emph{composed} coverage certificate (Theorem~\ref{thm:composed}) for
the joint safety event, with terms that map one-to-one onto the three modules: the soft
sensor's miscoverage, the prognostic bound's miscoverage, a similitude-departure penalty
contributed by each module's regime shift, and a selection penalty.
\item A resolution of the feedback objection. Under the causal timing that any physically
realizable loop satisfies, in which the decision is formed before the measurement it is
conditioned against, the selection penalty is exactly zero. The conformal-selection effect
is named and bounded rather than ignored, and it is shown to be active only under a
non-causal co-selection that real loops do not exhibit.
\item A conversion of the certificate into a single convex constraint, a \emph{$\psi$-budget}
on the similitude departure, that the modifier-adaptation RTO enforces directly
(Corollary~\ref{cor:floor}). The optimizer is barred from pushing economics hard enough to
invalidate its own sensor and life guarantees, and the flagship certificate and the
operating floor reduce to the same constraint.
\item A fully reproducible closed-loop demonstration on a calibrated debutanizer twin
running the real soft-sensor, prognostic, and RTO implementations. The certified floor
$1-(\alpha_1+\alpha_2)-\varepsilon = 0.75$ is met (joint safety coverage $0.988$) when the
optimizer enforces the budget and collapses (to $0.000$) when it does not, with the
contrast driven entirely by equipment life (Figure~\ref{fig:coverage}).
\end{enumerate}

The claim is sharp: none of the three guarantees implies system safety on its own; the
composition is the certificate, and a single convex budget makes it something the optimizer
must obey. The remainder of the paper reviews the relevant
conformal, soft-sensing, prognostic, and optimization literatures and positions the
contribution; develops the framework, the certificate, and the budget; describes the
debutanizer twin and the two-arm protocol; reports the closed-loop coverage result; and
discusses conservatism, the multi-cycle horizon, and scope.

\section{Related work}
\label{sec:related}

\paragraph{Conformal prediction beyond exchangeability.}
These methods give distribution-free, finite-sample coverage under the single
assumption that calibration and test data are exchangeable~\citep{vovk2005,shafer2008,
angelopoulos2023}, and regression variants such as conformalized quantile
regression~\citep{romano2019} and split-conformal methods~\citep{lei2018} make the
intervals locally adaptive. Exchangeability is the assumption that fails in feedback.
Two lines of work relax it. \citet{tibshirani2019} reweight the calibration scores when
the covariate shift is known, while \citet{barber2023} bound the coverage loss under an
arbitrary shift by a total-variation-type distance between the calibration and test
distributions. The per-module penalty $2L_j\|\Delta\psi_{j,k}\|$ used here is a
physically computable instance of that bound, with the abstract distributional distance
replaced by a similitude departure that is known in closed form from the operating
point. Adaptive conformal inference~\citep{gibbs2021} takes a complementary route and
guarantees long-run coverage online under arbitrary dependence by adjusting the
conformal quantile after each observation; we return to it for the multi-cycle horizon
in the discussion.

\paragraph{Soft sensing with calibrated uncertainty.}
Inferential (soft) sensors are standard practice for online estimation of qualities
that resist direct measurement~\citep{fortuna2007}, and the debutanizer is among the
most studied benchmarks, including recent machine-learning soft sensors for
product-quality monitoring on cyber-physical columns~\citep{rozanec2022}. The
constituent soft-sensor module~\citep{busico_m1} adds physics-anchored features and an
adaptive conformal calibration that holds coverage under delayed labels. We consume its
interval as the quality guarantee in the composition.

\paragraph{Data-driven prognostics.}
Statistical and data-driven estimation of remaining useful life is surveyed
in~\citet{si2011} and, for rotating machinery, in~\citet{lei2018machinery}. The
constituent prognostic module~\citep{busico_m2} contributes a similarity-calibrated
conformal lower bound on remaining life, which we consume as the equipment-life
guarantee.

\paragraph{Real-time optimization and health-aware control.}
RTO adjusts setpoints to track the economic optimum despite plant-model mismatch.
Modifier adaptation~\citep{marchetti2009,chachuat2009} corrects the optimizer with
measured plant gradients and is the method the constituent RTO
module~\citep{busico_m3} builds on. A broader lineage embeds economic objectives
directly in the feedback layer~\citep{skogestad2000}, and health- or constraint-aware
variants restrict operation to protect equipment. What this body of work does not yet
offer is a distribution-free coverage guarantee on the joint quality-and-life event that
\emph{composes} the uncertainty of the upstream estimators. The optimizer is made
health-aware, but the soft-sensor and prognostic guarantees are never propagated into a
certificate on the closed-loop outcome.

\paragraph{The gap.}
Across these literatures the per-layer guarantees exist and the health-aware optimizer
exists, yet their composition under feedback has not been addressed: a single per-cycle
certificate on the system-safety event, derived rather than asserted and enforceable as
one optimizer constraint. Supplying that composition is the contribution of this paper.

\section{Framework}
\label{sec:framework}

This section develops the central result: a per-cycle coverage certificate for
the closed loop whose terms map one-to-one onto the three constituent modules,
together with the constraint that renders it enforceable inside a real-time
optimizer. We first fix the loop and the event we certify, introduce the
dimensionless coordinate in which the optimizer's decisions move the modules'
operating regime, state the per-module conditional coverage these regimes
inherit, resolve the objection that feedback voids conformal validity, compose
the modules into a single certificate, and convert the certificate into an
optimizer constraint.

\subsection{The loop and the certified event}
\label{sec:loop}

\begin{figure}[t]
\centering
\includegraphics[width=\textwidth]{loop_schematic.png}
\caption{The integrated closed loop. The optimizer (Module~3) sets the decision
$u_k$; the plant returns a tray temperature to the soft sensor (Module~1) and a pump
load to the prognostic module (Module~2), whose interval and remaining-life bound feed
back to the optimizer. The reflux lever moves both modules' operating regimes at once
($\psi_1$ separation, $\psi_2$ pump load), and the composed certificate bounds the
coverage of the joint safety event $S_k$ in terms of the two regime departures.}
\label{fig:loop}
\end{figure}

At control cycle $k$ the real-time optimizer (RTO) sets the decision
$u_k = (R_k, D_k)$, the reflux ratio and distillate rate of a binary-key
distillation column. The plant returns tray temperatures $T_k$ and a true but
\emph{unmeasured} product quality $x_k$, the light-key (C4) mole fraction in the
bottoms. The reflux pump carries a true but \emph{unobserved} remaining useful
life $\rho_k$.

Three modules act on this loop (Figure~\ref{fig:loop}). The soft sensor (Module~1,~\citet{busico_m1})
issues a conformal interval $\hat C_1(T_k) = [\hat x_k - w_k,\, \hat x_k + w_k]$
for the quality with nominal miscoverage $\alpha_1$, so that marginally
$\mathbb{P}(x_k \in \hat C_1(T_k)) \ge 1-\alpha_1$. The prognostic module
(Module~2,~\citet{busico_m2}) issues a one-sided lower bound $\rho_k^{\mathrm{lo}}$
with nominal miscoverage $\alpha_2$, so that
$\mathbb{P}(\rho_k \ge \rho_k^{\mathrm{lo}}) \ge 1-\alpha_2$. The optimizer
(Module~3,~\citet{busico_m3}), a modifier-adaptation
RTO~\citep{marchetti2009,chachuat2009}, computes the next decision while enforcing
two \emph{surrogate} constraints stated on the estimates:
\begin{align}
\hat x_k + w_k &\le x^{\mathrm{spec}}, \label{eq:quality-surrogate}\\
\rho_k^{\mathrm{lo}} &\ge \rho_{\min}, \label{eq:health-surrogate}
\end{align}
where $x^{\mathrm{spec}}$ is the product specification and $\rho_{\min}$ the
remaining life required to reach the next maintenance window. Equation
\eqref{eq:quality-surrogate} backs the spec off by the interval half-width;
\eqref{eq:health-surrogate} is a health floor.

The quantity that actually matters for safe operation is neither estimate but the
joint event on the \emph{truths}: the product is in specification and the pump
survives to maintenance. We call this the system-safety event,
\begin{equation}
S_k \;=\; \{\, x_k \le x^{\mathrm{spec}} \,\} \;\cap\; \{\, \rho_k \ge \rho_{\min} \,\}.
\label{eq:safety-event}
\end{equation}
The surrogate constraints \eqref{eq:quality-surrogate}--\eqref{eq:health-surrogate}
act on estimates; safety \eqref{eq:safety-event} is about truths. The bridge from
one to the other is coverage, and the contribution of this paper is to make that
bridge a single, computable, enforceable guarantee.

\subsection{The similitude coordinate and why the optimizer controls it}
\label{sec:psi}

The modules' calibration is regime-dependent: a conformal interval calibrated at
one operating point need not retain its coverage at another. We track the
operating point in a dimensionless coordinate $\psi$ spanned by three standard
distillation quantities, each verified against Perry's Handbook~\citep{perrys}:
the relative volatility $\alpha(T)$ (Eq.~13-33), the Gilliland reflux coordinate
$\Psi_G = (R - R_{\min})/(R+1)$ (Eq.~13-30, Fig.~13-29), and the stripping factor
$S = KV/L$ (Eq.~13-44). Let $\psi_1\!\star$ and $\psi_2\!\star$ be the anchor
regimes at which Module~1 and Module~2 were calibrated, and define the residual
similitude departures
\begin{equation}
\Delta\psi_{1,k} = \psi_1(u_k) - \psi_1\!\star, \qquad
\Delta\psi_{2,k} = \psi_2(u_k) - \psi_2\!\star,
\label{eq:departure}
\end{equation}
the part of the regime shift not already absorbed by each module's internal
normalizer.

The key structural fact is that \emph{both} departures are functions of the
optimizer's decision. Raising the reflux ratio $R = L/D$ raises the reflux flow
$L$, which raises the reflux-pump duty; the pump-affinity laws (Perry's
Table~10-13: $Q \propto N$, $H \propto N^2$, $\mathrm{BHP} \propto N^3$) map flow
to mechanical load, and Module~2's degradation model maps load to wear. Their
composition makes $\psi_2(u_k)$ a function of $u_k$. The separation departure
$\psi_1(u_k)$ is a function of $u_k$ by construction. Thus the same lever that
the optimizer pulls for economic gain (reflux) simultaneously moves both
modules away from the regimes where their guarantees were calibrated. This endogeneity is the source of the difficulty in the next subsection and the
source of the controllability in the $\psi$-budget constraint developed below.

\subsection{Per-module conditional coverage}
\label{sec:permodule}

Within a single cycle the operating regime is held fixed while the modules
predict. The similarity-calibrated conformal certificate of~\citet{busico_m3},
applied conditionally on the cycle's own realized departure, gives
\begin{align}
\mathbb{P}\!\left(x_k \in \hat C_1(T_k) \,\middle|\, \Delta\psi_{1,k}\right)
  &\ge 1-\alpha_1 - 2L_1\|\Delta\psi_{1,k}\|, \label{eq:m1-cond}\\
\mathbb{P}\!\left(\rho_k \ge \rho_k^{\mathrm{lo}} \,\middle|\, \Delta\psi_{2,k}\right)
  &\ge 1-\alpha_2 - 2L_2\|\Delta\psi_{2,k}\|, \label{eq:m2-cond}
\end{align}
where $L_j$ is the Lipschitz constant of the score-quantile map in $\psi$-space,
estimated a~priori from a calibration sweep. The penalty $2L_j\|\Delta\psi_{j,k}\|$
is a physically computable instance of the non-exchangeability coverage gap of
\citet{barber2023}: their abstract distributional distance between calibration and
test is replaced here by the similitude departure, an engineering quantity known
in closed form from the operating point. We write
$E_{1,k} = \{x_k \in \hat C_1(T_k)\}$ and
$E_{2,k} = \{\rho_k \ge \rho_k^{\mathrm{lo}}\}$ for the two coverage events.

\subsection{Why feedback does not break the certificate}
\label{sec:causal}

Because $u_{k+1}$ is computed from $(\hat x_k, \rho_k^{\mathrm{lo}})$, the departure
sequence $\{\Delta\psi_{j,k}\}$ is endogenous, dependent, and controlled, and the
cycle-$k$ test point is not exchangeable with the offline calibration set. This is
the standard reason conformal guarantees are doubted under feedback. The objection
is resolved in two parts.

First, \emph{conditioning}. The bounds \eqref{eq:m1-cond}--\eqref{eq:m2-cond} are
conditional on $\Delta\psi_{j,k}$ and are agnostic to the mechanism that produced
the regime. Endogeneity makes $\Delta\psi_{j,k}$ random; it does not invalidate a
statement already conditioned on it.

Second, \emph{causal timing}. We make explicit the timing assumption that the
closed loop must satisfy.
\begin{assumption}[Causal two-step timing]
\label{ass:causal}
The decision $u_k$ is $\sigma(\mathcal{F}_{k-1})$-measurable, a function of
information available strictly before cycle $k$, and the cycle-$k$ measurement
noise $\varepsilon_k$ satisfies $\varepsilon_k \perp \mathcal{F}_{k-1}$.
\end{assumption}
Under Assumption~\ref{ass:causal}, $u_k$ and $\varepsilon_k$ are conditionally
independent given $\Delta\psi_k$ (itself a function of $u_k$), so conditioning on
the regime leaves $\varepsilon_k$ at its nominal regime-conditional distribution:
\begin{equation}
u_k \perp \varepsilon_k \,\big|\, \Delta\psi_k \;\;\Longrightarrow\;\; \Delta_{\mathrm{sel},k} = 0,
\label{eq:no-selection}
\end{equation}
where $\Delta_{\mathrm{sel},k}$ is a selection penalty that is nonzero only in the
alternative configuration in which the optimizer co-selects $u_k$ on a quantity
computed from the \emph{same} sample used for the coverage check. This delineates
exactly when the conformal-selection effect of~\citet{busico_m3} is active: under
realistic causal plant timing it vanishes; the framework still names and bounds
it. Assumption~\ref{ass:causal} is not a hope about the implementation; it is a
property of the synchronous control loop used here, which reads $u_k$ from the
solution computed at the close of cycle $k-1$, applies it, and only then realizes
the cycle-$k$ measurement, so that the decision is constructed before the noise it
is conditioned against.

\subsection{The composed certificate}
\label{sec:theorem}

\begin{theorem}[Composed coverage certificate]
\label{thm:composed}
Suppose the optimizer enforces the surrogate constraints
\eqref{eq:quality-surrogate}--\eqref{eq:health-surrogate}, the per-module
conditional bounds \eqref{eq:m1-cond}--\eqref{eq:m2-cond} hold, and
Assumption~\ref{ass:causal} holds. Then the system-safety event
\eqref{eq:safety-event} satisfies, conditionally on the cycle's realized
departures,
\begin{equation}
\mathbb{P}\!\left(S_k \,\middle|\, \Delta\psi_{1,k}, \Delta\psi_{2,k}\right)
\;\ge\; 1 - (\alpha_1 + \alpha_2)
       - 2\!\left(L_1\|\Delta\psi_{1,k}\| + L_2\|\Delta\psi_{2,k}\|\right),
\label{eq:certificate-cond}
\end{equation}
and, marginalizing over the endogenous regime,
\begin{equation}
\mathbb{P}(S_k) \;\ge\; 1 - (\alpha_1 + \alpha_2)
       - 2\!\left(L_1\,\mathbb{E}\|\Delta\psi_{1,k}\| + L_2\,\mathbb{E}\|\Delta\psi_{2,k}\|\right).
\label{eq:certificate-marg}
\end{equation}
\end{theorem}

\begin{proof}
On the enforced surrogate, each coverage event implies its safety half. If
$E_{1,k}$ holds then $x_k \le \hat x_k + w_k \le x^{\mathrm{spec}}$ by
\eqref{eq:quality-surrogate}; if $E_{2,k}$ holds then
$\rho_k \ge \rho_k^{\mathrm{lo}} \ge \rho_{\min}$ by \eqref{eq:health-surrogate}.
Hence $E_{1,k}\cap E_{2,k} \subseteq S_k$, and so
$\mathbb{P}(S_k \mid \cdot) \ge \mathbb{P}(E_{1,k}\cap E_{2,k}\mid\cdot)$.
By the union bound applied to the complementary miscoverage events,
$\mathbb{P}(E_{1,k}\cap E_{2,k}\mid\cdot) \ge 1 - \mathbb{P}(E_{1,k}^c\mid\cdot)
- \mathbb{P}(E_{2,k}^c\mid\cdot)$. Substituting
\eqref{eq:m1-cond}--\eqref{eq:m2-cond}, and using
$\Delta_{\mathrm{sel},k}=0$ from \eqref{eq:no-selection} so the conditional
miscoverages are exactly those of the offline calibration, gives
\eqref{eq:certificate-cond}. Taking expectations over $\Delta\psi_{1,k},
\Delta\psi_{2,k}$ and applying Jensen's inequality to the (linear) departure terms
yields \eqref{eq:certificate-marg}.
\end{proof}

Equation \eqref{eq:certificate-cond} is the unification. Its terms map one-to-one
onto the three modules (Table~\ref{tab:terms}): the base coverage of $\hat C_1$ is
Module~1; the one-sided RUL bound and its miscoverage $\alpha_2$ are Module~2; the
similitude-departure penalty $2L_j\|\Delta\psi_{j,k}\|$ is the per-module,
per-cycle instance of the Module~3 certificate; and the (here vanishing) selection
penalty $\Delta_{\mathrm{sel},k}$ is the Module~3 conformal-selection effect, active
only when the optimizer co-selects on the test residual. No single module states
this inequality; the composition is the contribution.

\begin{table}[t]
\centering
\caption{Each term of the composed certificate \eqref{eq:certificate-cond}
originates in a distinct module.}
\label{tab:terms}
\begin{tabular}{ll}
\toprule
Term & Origin \\
\midrule
$\alpha_1$ (base coverage of $\hat C_1$) & Module~1, conformal soft sensor \\
$\alpha_2$ and the one-sided RUL bound & Module~2, calibrated prognostics \\
$2L_j\|\Delta\psi_{j,k}\|$ (similitude departure) & Module~3, SCC certificate (per cycle) \\
$\Delta_{\mathrm{sel},k}$ (selection penalty) & Module~3, selection effect ($=0$ here) \\
\bottomrule
\end{tabular}
\end{table}

\subsection{The $\psi$-budget constraint}
\label{sec:budget}

Because each departure $\Delta\psi_{j,k} = \psi_j(u_k) - \psi_j\!\star$ is a known
function of the decision $u_k$, the coverage penalty in
\eqref{eq:certificate-cond} is itself controllable. Imposing the single convex
constraint
\begin{equation}
2\!\left(L_1\|\psi_1(u) - \psi_1\!\star\| + L_2\|\psi_2(u) - \psi_2\!\star\|\right) \;\le\; \varepsilon
\label{eq:psi-budget}
\end{equation}
inside the modifier-adaptation NLP yields the following.

\begin{corollary}[Certified floor]
\label{cor:floor}
Under the hypotheses of Theorem~\ref{thm:composed}, if the optimizer additionally
enforces \eqref{eq:psi-budget}, then for every cycle
\begin{equation}
\mathbb{P}(S_k) \;\ge\; 1 - (\alpha_1 + \alpha_2) - \varepsilon.
\label{eq:floor}
\end{equation}
\end{corollary}

\begin{proof}
Constraint \eqref{eq:psi-budget} bounds the departure term in
\eqref{eq:certificate-cond} by $\varepsilon$ for every realized $u_k$, hence for
every realized $(\Delta\psi_{1,k},\Delta\psi_{2,k})$; the bound
\eqref{eq:floor} then holds conditionally and therefore marginally.
\end{proof}

Corollary~\ref{cor:floor} converts the certificate into an operating rule: the
optimizer maximizes economic profit \emph{subject to remaining inside the
similitude-departure budget that certifies its own coverage}. Formally, it may
not optimize so aggressively that it invalidates its own soft sensor and its own
RUL guarantee. The radius $\varepsilon$ is the chosen coverage budget, traded
directly against economics. Theorem~\ref{thm:composed} and
Corollary~\ref{cor:floor} share one implementation: the flagship per-cycle
guarantee and the operational floor are the same constraint
\eqref{eq:psi-budget} added to the NLP.

\section{Case study: a debutanizer digital twin}
\label{sec:case}

We demonstrate the certificate of Theorem~\ref{thm:composed} and the operating
rule of Corollary~\ref{cor:floor} in closed loop on a calibrated digital twin of a
debutanizer column. The three modules are the real implementations of the
constituent papers; the plant is a twin rather than an operating column. This is
a deliberate scoping choice, justified at the end of this section: the certificate
is a statement about how three coverage guarantees \emph{compose} in feedback, and
a twin that responds to the optimizer's setpoint moves is the appropriate
instrument for that question, whereas component accuracy is established on
operational data in the constituent papers.

\subsection{The plant twin}
\label{sec:twin}

The plant truth is a Fenske--Underwood--Gilliland (FUG) shortcut column model
($\alpha_{\mathrm{LK}}=5.97$, eight equilibrium stages, feed $100$~kmol\,h$^{-1}$,
light-key feed fraction $0.35$): for a decision $(R,D)$ it returns the bottoms
light-key fraction $x$, which is the unmeasured quality the soft sensor must
infer. The inferential measurement is a tray-6 temperature, modelled as cooling
monotonically with reflux as the split sharpens and landing in the
$[100.5, 111.5]\,^\circ$C window of the soft sensor's physics bridge; it carries a
$0.25\,^\circ$C sensor noise, and the bottoms truth carries a
$1.0\times10^{-3}$ process noise. The thermophysical properties (relative
volatility, heavy-key $K$-value, and liquid viscosity) are evaluated from
DIPPR-101 vapor pressures for the $n$-butane/$n$-hexane key pair.

The reflux pump degrades as a function of the optimizer's decision, closing in the
twin the endogeneity loop identified in the framework. Damage accrues at a rate
proportional to $\mathrm{load}^{p}$, where the load is the affinity brake-horsepower
ratio $\mathrm{load} = (Q/Q_{\mathrm{ref}})^{3}$ (Perry's Table~10-13) and the
exponent $p=2$ follows from rolling-element bearing fatigue (the $L_{10}$ life
varies as $\mathrm{load}^{-3}$) under a hydraulic radial load that scales as $Q^{2}$
by the affinity head law, so that damage scales as $Q^{6}$ in reflux flow. The base
rate is set to a fresh-nominal remaining life of $\approx\!1000$~h; these
coefficients are calibrated to the bearing-degradation timescale of the prognostic
module. Degradation severity drives a five-feature synthesizer whose output the
prognostic module scores.

\subsection{The three modules on the twin}
\label{sec:modules}

\emph{Module~1 (soft sensor).} The real physics-anchored soft
sensor~\citep{busico_m1} is trained on twin samples: a standardized ridge estimator
on four physics-derived features (the C4 bubble-point fraction, the relative
volatility, the stripping factor, and the delayed normalized tray temperature),
trained over $700$ operating points spanning the box and split $450$ for fitting
and $250$ for the conformal calibration residuals. On held-out twin data its adaptive conformal
interval attains empirical coverage $0.898$ against a $0.90$ target, with a mean
half-width of $7.1\times10^{-3}$, comparable to the product specification, so the
interval materially backs off the optimizer through
\eqref{eq:quality-surrogate}.

\emph{Module~2 (prognostics).} The real predictive-maintenance
stack~\citep{busico_m2} is built from its canonical configuration: a five-feature
Hotelling $T^{2}$ health index fitted on a healthy reference (alarm threshold
$T^{2}=15.1$) and a trajectory-similarity remaining-life estimator over a library of
run-to-failure arcs. Driven by the synthesizer, the health index rises
monotonically with severity and alarms at severity $\approx\!0.71$.

\emph{Module~3 (optimizer).} The real modifier-adaptation
RTO~\citep{busico_m3} optimizes economic profit over the FUG response surface
(a log-quadratic fit in $(R,D)$, coefficient of determination $0.982$), enforcing
the backed-off quality constraint \eqref{eq:quality-surrogate} and, in the
health-constrained arm, the $\psi$-budget \eqref{eq:psi-budget}. The program is
solved to local optimality at each cycle.

\subsection{Protocol}
\label{sec:protocol}

The demonstration contrasts two arms that are identical in every respect except the
$\psi$-budget. In the \emph{health-blind} arm the budget is inactive
($\varepsilon\to\infty$) and the optimizer maximizes economics subject only to the
quality surrogate. In the \emph{health-constrained} arm the budget is active with
radius $\varepsilon=0.05$. Each arm is run for $8$ random seeds of $80$ control
cycles; the coverage harness records, per cycle, whether the system-safety event
\eqref{eq:safety-event} holds against the truths, and aggregates the empirical
$S_k$-coverage with a Wilson interval. Both arms share the parameters of
Table~\ref{tab:params}. With $\alpha_1=\alpha_2=0.1$ and $\varepsilon=0.05$, the
certified floor of Corollary~\ref{cor:floor} is
$1-(\alpha_1+\alpha_2)-\varepsilon = 0.75$.

The Lipschitz constant of the soft-sensor score-quantile in $\psi$-space, estimated
a~priori from a regime sweep across the operating box, is
$L_1 \approx 0.025$ (Figure~\ref{fig:l1}). The demonstration uses a deliberately conservative common value
$L_1=L_2=0.15$ in the enforced budget \eqref{eq:psi-budget}; because this exceeds
the swept estimate, the budget is tighter than the data require and the certified
floor holds a fortiori.

\begin{figure}[t]
\centering
\includegraphics[width=0.82\textwidth]{l1_sweep.png}
\caption{A-priori estimate of the soft-sensor Lipschitz constant $L_1$. Each point is
an operating regime, placed at its similitude departure $\|\Delta\psi_1\|$ from the
calibration anchor (star); the vertical coordinate is the conformal score quantile at
level $1-\alpha_1$. The quantile is nearly flat in $\|\Delta\psi_1\|$, and every regime
falls inside the Lipschitz cone (shaded), whose slope $L_1\approx0.025$ is set by the
binding pair (red). The enforced budget \eqref{eq:psi-budget} uses $L_1=L_2=0.15$,
roughly six times this estimate, so the coverage penalty it charges over-bounds the
true loss.}
\label{fig:l1}
\end{figure}

\begin{table}[t]
\centering
\caption{Demonstration parameters. Both arms are identical except the
$\psi$-budget radius $\varepsilon$.}
\label{tab:params}
\begin{tabular}{lll}
\toprule
Symbol & Meaning & Value \\
\midrule
$\alpha_1$ & soft-sensor miscoverage & $0.10$ \\
$\alpha_2$ & prognostic miscoverage & $0.10$ \\
$\varepsilon$ & $\psi$-budget radius (constrained arm) & $0.05$ \\
$x^{\mathrm{spec}}$ & bottoms C4 specification & $8.0\times10^{-3}$ \\
$\rho_{\min}$ & remaining-life floor & $200$~h \\
$L_1=L_2$ & Lipschitz constants used (conservative) & $0.15$ \\
$(R,D)_{\mathrm{nom}}$ & nominal optimum & $(2.037,\,33.49)$ \\
seeds $\times$ cycles & Monte-Carlo replicates & $8\times80$ \\
\bottomrule
\end{tabular}
\end{table}

\subsection{Scope of the demonstration}
\label{sec:scope}

The closed-loop demonstration runs on a calibrated digital twin; the
\emph{accuracy} of each module is validated on operational data in the constituent
papers~\citep{busico_m1,busico_m2,busico_m3}. Two reasons make the twin the correct
instrument here. First, a closed-loop optimizer requires a plant that
\emph{responds} to its setpoint moves, because the endogeneity identified in the
framework cannot be exercised against a frozen historical dataset. Second, the
certificate concerns the \emph{composition} of three coverage guarantees under
feedback, a property of the loop architecture rather than of any module's point
accuracy; the twin isolates exactly that property while remaining fully
reproducible from open code with no proprietary data. Physical realization on a
dynamic high-fidelity simulator and on plant data is the subject of future work.

\section{Results}
\label{sec:results}

The central finding is that the certified floor of Corollary~\ref{cor:floor} is
honoured empirically when the optimizer enforces the $\psi$-budget and is violated
when it does not. Table~\ref{tab:coverage} reports the two arms against the floor;
Figure~\ref{fig:coverage} shows the coverage decomposition and the per-cycle
mechanism.

\begin{table}[t]
\centering
\caption{Closed-loop coverage over $8$ seeds $\times\,80$ cycles ($N=640$ per arm).
The health-constrained arm meets the certified floor $0.75$; the health-blind arm
does not. Brackets are Wilson $95\%$ intervals.}
\label{tab:coverage}
\begin{tabular}{lccc}
\toprule
 & Health-blind & Health-constrained & Certified floor \\
 & ($\varepsilon\to\infty$) & ($\varepsilon=0.05$) & \\
\midrule
$S_k$ coverage & $0.000\,[0.000,0.006]$ & $\mathbf{0.988\,[0.976,0.994]}$ & $0.750$ \\
\quad meets floor? & no & \textbf{yes} & \\
Quality coverage $\{x_k\le x^{\mathrm{spec}}\}$ & $0.850$ & $0.988$ & \\
RUL coverage $\{\rho_k\ge\rho_{\min}\}$ & $0.013$ & $1.000$ & \\
$\psi$-budget binding & $0\%$ & $100\%$ & \\
Mean final severity & $0.700$ & $0.097$ & \\
\bottomrule
\end{tabular}
\end{table}

\begin{figure}[t]
\centering
\includegraphics[width=\textwidth]{twin_coverage.png}
\caption{Composed coverage certificate on the debutanizer twin.
(a) Empirical coverage of the joint safety event $S_k$ and its two halves for the
two arms, against the certified floor $1-(\alpha_1+\alpha_2)-\varepsilon=0.75$
(dashed). (b) True reflux-pump remaining life per cycle; the health-blind arm falls
below the floor $\rho_{\min}=200$~h within two cycles, the health-constrained arm
holds $700$--$1000$~h. (c) Degradation severity per cycle.}
\label{fig:coverage}
\end{figure}

\subsection{The certificate holds under the budget}
\label{sec:res-cert}

Under the $\psi$-budget the joint safety event holds on $0.988$ of cycles
($95\%$ Wilson interval $[0.976, 0.994]$), comfortably above the certified floor of
$0.75$. Both halves of $S_k$ are satisfied simultaneously: quality coverage is
$0.988$ and remaining-life coverage is $1.000$. The budget is active on every
constrained cycle, holding the operating point on the certified reference; the
reflux pump ends the horizon at a mean severity of $0.097$, far from its alarm.

The empirical coverage exceeds the floor by a wide margin. This is expected and
consistent with the construction: the certificate is a valid lower bound, and two
sources of conservatism inflate the gap: the Bonferroni union over the two
miscoverage events (Theorem~\ref{thm:composed}), which is tight only when the two
failures are disjoint, and the deliberately conservative Lipschitz constants
($L_1=L_2=0.15$ against a swept $L_1\approx0.025$). What matters in
Table~\ref{tab:coverage} is therefore the verdict, that the floor is met with a
margin, not the size of that margin.

\subsection{The mechanism: why the blind optimizer fails}
\label{sec:res-mech}

The health-blind arm exposes what the budget protects against. Facing the same
tight specification, the unconstrained optimizer drives the reflux ratio to its
upper bound ($R\to3.0$, reflux flow $\approx\!97$~kmol\,h$^{-1}$) to maximize C4
recovery. By the affinity and bearing-fatigue scaling of the pump model, the
pump load rises steeply: severity climbs to $0.70$ and the true remaining life
collapses below the $200$~h floor within two cycles
(Figure~\ref{fig:coverage}b,c). The remaining-life half of $S_k$ therefore fails on
$98.7\%$ of cycles, and the joint event never holds: $S_k$ coverage is exactly
$0.000$, with no certificate available because the floor $1-(\alpha_1+\alpha_2)-
\varepsilon$ is meaningless once $\varepsilon$ is unbounded. Notably the blind arm's
quality coverage ($0.850$) is itself \emph{below} the soft sensor's nominal
$1-\alpha_1=0.90$: pushed to the edge of the operating box, the optimizer drives the
soft sensor outside the regime where its interval was calibrated, which is precisely the
non-exchangeability penalty $2L_1\|\Delta\psi_1\|$ that the budget exists to bound.

The contrast is the whole point. The two arms differ only in whether the optimizer
is allowed to leave the similitude-departure budget; that single convex constraint
moves the joint safety coverage from $0.000$ to $0.988$, and it does so by
preserving \emph{both} guarantees at once, quality \emph{and} remaining life,
rather than trading one for the other. The flagship theorem and the operating floor
are, as Corollary~\ref{cor:floor} states, the same constraint.

\section{Discussion}
\label{sec:discussion}

\paragraph{The certificate is a conservative lower bound.}
The empirical coverage under the budget ($0.988$) sits well above the certified floor
($0.75$), and the construction explains the gap. Theorem~\ref{thm:composed} composes the
two miscoverage events through a union bound, which is tight only when the two failures
are disjoint. In practice the soft-sensor and prognostic failures are positively
correlated in the regimes that matter, because aggressive operation pushes the soft
sensor out of its calibration regime and loads the pump at the same time, so the two
modules tend to be stressed together. Positive failure correlation makes the true joint
coverage exceed the Bonferroni bound, and the deliberately conservative Lipschitz
constants ($L_1=L_2=0.15$ against a swept $L_1\approx0.025$) widen the gap further. The
certificate is therefore a guarantee, a floor that is provably met, rather than a
prediction of the coverage actually achieved. Tightening it is a concrete next step: a
joint failure-correlation model, built from a copula on the two nonconformity scores,
from multivariate conformal prediction, or from conformal risk control under dependence,
would replace the union bound with a dependence-aware composition and recover much of the
margin.

\paragraph{From per-cycle to the operating horizon.}
Theorem~\ref{thm:composed} is a per-cycle, marginal guarantee. A campaign runs over many
cycles, and a naive union of the per-cycle bound over a horizon of $K$ cycles degrades
linearly in $K$. The route to a genuine horizon guarantee is adaptive conformal
inference~\citep{gibbs2021}, which the soft-sensor module already
implements~\citep{busico_m1}. By adjusting the conformal quantile online it secures
long-run coverage under arbitrary, unspecified dependence, which is exactly the
dependence that feedback induces across cycles. Lifting the per-cycle certificate to a
horizon statement under this mechanism, in a dynamic rather than quasi-static setting, is
the principal direction of ongoing work, and it is what connects the present result to a
physically deployed implementation.

\paragraph{The selection term, and when it is not zero.}
The selection penalty $\Delta_{\mathrm{sel},k}$ vanishes here as a property of the loop
(Assumption~\ref{ass:causal}), not as an idealization, because the synchronous controller
forms each decision from the previous cycle's information and applies it before the
current measurement is realized. Carrying the term explicitly is worthwhile because it
identifies precisely when conformal validity would instead be lost, namely under a
non-causal co-selection in which the optimizer chooses the decision using the same
residual that the coverage check then evaluates. Most treatments of conformal prediction
in feedback either assume the difficulty away or abandon the guarantee. Naming the
selection effect, bounding it, and showing that it is null under realizable timing is what
makes the closed-loop certificate honest.

\paragraph{Scope and generality.}
The certificate is instantiated for binary-key distillation through a similitude
coordinate built from the relative volatility, the Gilliland reflux coordinate, and the
stripping factor. That coordinate is the unit-specific ingredient. The composition logic
itself, namely per-module conditional coverage, the causal-timing argument, the Bonferroni
composition, and the convex $\psi$-budget, is unit-agnostic. Extending the framework to
multicomponent separations or to other unit operations requires identifying the
appropriate dimensionless groups for those units, after which the same machinery applies.
A second refinement concerns the Lipschitz constants, estimated here a~priori from a
regime sweep; an online estimator that re-identifies them as the plant ages would let the
budget adapt, and it pairs naturally with the adaptive-conformal horizon extension above.
The budget radius $\varepsilon$ is, finally, an explicit economic dial. It trades the
certified coverage level directly against the profit the optimizer is permitted to chase,
and exposing that trade-off quantitatively is one of the practical contributions of the
framework.

\section{Conclusion}
\label{sec:conclusion}

Soft sensing, prognostics, and real-time optimization each ship a calibrated uncertainty
guarantee, but a closed-loop operator needs a guarantee on the joint system-safety event,
namely an in-specification product together with surviving equipment, and that joint
guarantee does not follow from the three taken separately. This paper supplied it. The
per-cycle composed coverage certificate of Theorem~\ref{thm:composed} states the joint
guarantee explicitly, with terms that map one-to-one onto the three modules, and a
causal-timing argument shows that feedback does not void it: the conformal-selection
penalty is exactly zero whenever the decision is formed before the measurement it is
conditioned against. Corollary~\ref{cor:floor} then converts the certificate into a
single convex $\psi$-budget that a modifier-adaptation optimizer enforces directly, so
that the flagship guarantee and the operating floor are one and the same constraint, and
the optimizer is barred from pushing economics hard enough to invalidate its own sensor
and life guarantees.

On a calibrated debutanizer twin running the real soft-sensor, prognostic, and
optimization implementations, the certified floor $1-(\alpha_1+\alpha_2)-\varepsilon =
0.75$ was met empirically, with joint safety coverage of $0.988$ under the budget and
$0.000$ without it, the entire contrast carried by reflux-pump life. None of the three
guarantees implies system safety on its own; the composition is the certificate.

Three directions follow. The per-cycle certificate should be lifted to a horizon
guarantee under adaptive conformal inference and realized on a dynamic high-fidelity
simulator and on plant data, moving from the quasi-static demonstration given here to a
physically deployed loop. The Bonferroni composition should be tightened with a
dependence-aware model of the joint soft-sensor and prognostic failure, which would
recover the margin that positive failure correlation currently leaves unused. Finally,
the framework should be carried from a single column to a plant, with the coverage budget
allocated across units. That last step is what turns a certificate for one column into a
certificate for an integrated process.


\section*{Data and code availability}
The reference implementation of all three modules, the closed-loop integration
layer, the coverage harness, and the single script that regenerates every number
and figure in this paper are available at \url{github.com/beebzy-droid/IPIS}.

\section*{Declaration of competing interest}
The author declares no competing financial interests or personal relationships that
could have influenced the work reported in this paper.

\begin{thebibliography}{19}
\expandafter\ifx\csname natexlab\endcsname\relax\def\natexlab#1{#1}\fi
\providecommand{\url}[1]{\texttt{#1}}
\providecommand{\href}[2]{#2}
\providecommand{\path}[1]{#1}
\providecommand{\DOIprefix}{doi:}
\providecommand{\ArXivprefix}{arXiv:}
\providecommand{\URLprefix}{URL: }
\providecommand{\Pubmedprefix}{pmid:}
\providecommand{\doi}[1]{\href{http://dx.doi.org/#1}{\path{#1}}}
\providecommand{\Pubmed}[1]{\href{pmid:#1}{\path{#1}}}
\providecommand{\bibinfo}[2]{#2}
\ifx\xfnm\relax \def\xfnm[#1]{\unskip,\space#1}\fi
%Type = Book
\bibitem[{Vovk et~al.(2005)Vovk, Gammerman, and Shafer}]{vovk2005}
\bibinfo{author}{V.~Vovk}, \bibinfo{author}{A.~Gammerman},
  \bibinfo{author}{G.~Shafer}, \bibinfo{title}{Algorithmic Learning in a Random
  World}, \bibinfo{publisher}{Springer}, \bibinfo{address}{New York},
  \bibinfo{year}{2005}.
%Type = Article
\bibitem[{Shafer and Vovk(2008)}]{shafer2008}
\bibinfo{author}{G.~Shafer}, \bibinfo{author}{V.~Vovk},
\newblock \bibinfo{title}{A tutorial on conformal prediction},
\newblock \bibinfo{journal}{Journal of Machine Learning Research}
  \bibinfo{volume}{9} (\bibinfo{year}{2008}) \bibinfo{pages}{371--421}.
%Type = Article
\bibitem[{Angelopoulos and Bates(2023)}]{angelopoulos2023}
\bibinfo{author}{A.~N. Angelopoulos}, \bibinfo{author}{S.~Bates},
\newblock \bibinfo{title}{Conformal prediction: A gentle introduction},
\newblock \bibinfo{journal}{Foundations and Trends in Machine Learning}
  \bibinfo{volume}{16} (\bibinfo{year}{2023}) \bibinfo{pages}{494--591}.
%Type = Article
\bibitem[{Romano et~al.(2019)Romano, Patterson, and Cand\`es}]{romano2019}
\bibinfo{author}{Y.~Romano}, \bibinfo{author}{E.~Patterson},
  \bibinfo{author}{E.~J. Cand\`es},
\newblock \bibinfo{title}{Conformalized quantile regression},
\newblock \bibinfo{journal}{Advances in Neural Information Processing Systems}
  \bibinfo{volume}{32} (\bibinfo{year}{2019}).
%Type = Article
\bibitem[{Lei et~al.(2018)Lei, G'Sell, Rinaldo, Tibshirani, and
  Wasserman}]{lei2018}
\bibinfo{author}{J.~Lei}, \bibinfo{author}{M.~G'Sell},
  \bibinfo{author}{A.~Rinaldo}, \bibinfo{author}{R.~J. Tibshirani},
  \bibinfo{author}{L.~Wasserman},
\newblock \bibinfo{title}{Distribution-free predictive inference for
  regression},
\newblock \bibinfo{journal}{Journal of the American Statistical Association}
  \bibinfo{volume}{113} (\bibinfo{year}{2018}) \bibinfo{pages}{1094--1111}.
%Type = Article
\bibitem[{Tibshirani et~al.(2019)Tibshirani, Foygel~Barber, Cand\`es, and
  Ramdas}]{tibshirani2019}
\bibinfo{author}{R.~J. Tibshirani}, \bibinfo{author}{R.~Foygel~Barber},
  \bibinfo{author}{E.~Cand\`es}, \bibinfo{author}{A.~Ramdas},
\newblock \bibinfo{title}{Conformal prediction under covariate shift},
\newblock \bibinfo{journal}{Advances in Neural Information Processing Systems}
  \bibinfo{volume}{32} (\bibinfo{year}{2019}).
%Type = Article
\bibitem[{Barber et~al.(2023)Barber, Cand\`es, Ramdas, and
  Tibshirani}]{barber2023}
\bibinfo{author}{R.~F. Barber}, \bibinfo{author}{E.~J. Cand\`es},
  \bibinfo{author}{A.~Ramdas}, \bibinfo{author}{R.~J. Tibshirani},
\newblock \bibinfo{title}{Conformal prediction beyond exchangeability},
\newblock \bibinfo{journal}{The Annals of Statistics} \bibinfo{volume}{51}
  (\bibinfo{year}{2023}) \bibinfo{pages}{816--845}.
%Type = Inproceedings
\bibitem[{Gibbs and Cand\`es(2021)}]{gibbs2021}
\bibinfo{author}{I.~Gibbs}, \bibinfo{author}{E.~J. Cand\`es},
\newblock \bibinfo{title}{Adaptive conformal inference under distribution
  shift},
\newblock in: \bibinfo{booktitle}{Advances in Neural Information Processing
  Systems}, volume~\bibinfo{volume}{34}, \bibinfo{year}{2021}, pp.
  \bibinfo{pages}{1660--1672}.
%Type = Article
\bibitem[{Fortuna et~al.(2007)Fortuna, Graziani, Rizzo, and
  Xibilia}]{fortuna2007}
\bibinfo{author}{L.~Fortuna}, \bibinfo{author}{S.~Graziani},
  \bibinfo{author}{A.~Rizzo}, \bibinfo{author}{M.~G. Xibilia},
\newblock \bibinfo{title}{Soft sensors for monitoring and control of industrial
  processes}  (\bibinfo{year}{2007}).
%Type = Article
\bibitem[{Ro{\v{z}}anec et~al.(2022)}]{rozanec2022}
\bibinfo{author}{J.~M. Ro{\v{z}}anec}, et~al.,
\newblock \bibinfo{title}{Cyber-physical {LPG} debutanizer distillation
  columns: machine-learning-based soft sensors for product quality monitoring},
\newblock \bibinfo{journal}{Applied Sciences}  (\bibinfo{year}{2022}).
%Type = Misc
\bibitem[{Busico(2026)}]{busico_m1}
\bibinfo{author}{B.~Busico}, \bibinfo{title}{{When does a calibrated soft
  sensor keep its promise? A negative-control study of validity without
  accuracy under drift and delayed labels}}, \bibinfo{year}{2026}.
  \bibinfo{note}{Manuscript JPROCONT-D-26-00618, Journal of Process Control
  (transfer from CACE-D-26-00944); reframed as a negative-control study}.
%Type = Article
\bibitem[{Si et~al.(2011)Si, Wang, Hu, and Zhou}]{si2011}
\bibinfo{author}{X.-S. Si}, \bibinfo{author}{W.~Wang}, \bibinfo{author}{C.-H.
  Hu}, \bibinfo{author}{D.-H. Zhou},
\newblock \bibinfo{title}{Remaining useful life estimation--a review on the
  statistical data driven approaches},
\newblock \bibinfo{journal}{European Journal of Operational Research}
  \bibinfo{volume}{213} (\bibinfo{year}{2011}) \bibinfo{pages}{1--14}.
%Type = Article
\bibitem[{Lei et~al.(2018)Lei, Li, Guo, Li, Yan, and Lin}]{lei2018machinery}
\bibinfo{author}{Y.~Lei}, \bibinfo{author}{N.~Li}, \bibinfo{author}{L.~Guo},
  \bibinfo{author}{N.~Li}, \bibinfo{author}{T.~Yan}, \bibinfo{author}{J.~Lin},
\newblock \bibinfo{title}{Machinery health prognostics: A systematic review
  from data acquisition to rul prediction},
\newblock \bibinfo{journal}{Mechanical Systems and Signal Processing}
  \bibinfo{volume}{104} (\bibinfo{year}{2018}) \bibinfo{pages}{799--834}.
%Type = Misc
\bibitem[{Busico(2026)}]{busico_m2}
\bibinfo{author}{B.~Busico}, \bibinfo{title}{{Similarity-Calibrated Conformal
  prediction: data-free coverage guarantees for remaining-useful-life intervals
  under operating-regime transfer}}, \bibinfo{year}{2026}.
  \bibinfo{note}{Manuscript JRESS-D-26-04700, submitted to Reliability
  Engineering \& System Safety}.
%Type = Article
\bibitem[{Marchetti et~al.(2009)Marchetti, Chachuat, and
  Bonvin}]{marchetti2009}
\bibinfo{author}{A.~Marchetti}, \bibinfo{author}{B.~Chachuat},
  \bibinfo{author}{D.~Bonvin},
\newblock \bibinfo{title}{Modifier-adaptation methodology for real-time
  optimization},
\newblock \bibinfo{journal}{Industrial \& Engineering Chemistry Research}
  \bibinfo{volume}{48} (\bibinfo{year}{2009}) \bibinfo{pages}{6022--6033}.
%Type = Article
\bibitem[{Chachuat et~al.(2009)Chachuat, Srinivasan, and Bonvin}]{chachuat2009}
\bibinfo{author}{B.~Chachuat}, \bibinfo{author}{B.~Srinivasan},
  \bibinfo{author}{D.~Bonvin},
\newblock \bibinfo{title}{Adaptation strategies for real-time optimization},
\newblock \bibinfo{journal}{Computers \& Chemical Engineering}
  \bibinfo{volume}{33} (\bibinfo{year}{2009}) \bibinfo{pages}{1557--1567}.
%Type = Misc
\bibitem[{Busico(2026)}]{busico_m3}
\bibinfo{author}{B.~Busico}, \bibinfo{title}{{Conditionally calibrated
  conformal back-offs for chance-constrained real-time optimisation under
  unmeasured disturbances}}, \bibinfo{year}{2026}. \bibinfo{note}{Manuscript
  CACE-D-26-01040, submitted to Computers \& Chemical Engineering}.
%Type = Article
\bibitem[{Skogestad(2000)}]{skogestad2000}
\bibinfo{author}{S.~Skogestad},
\newblock \bibinfo{title}{Plantwide control: the search for the self-optimizing
  control structure},
\newblock \bibinfo{journal}{Journal of Process Control} \bibinfo{volume}{10}
  (\bibinfo{year}{2000}) \bibinfo{pages}{487--507}.
%Type = Book
\bibitem[{Green and Southard(2019)}]{perrys}
\bibinfo{author}{D.~W. Green}, \bibinfo{author}{M.~Z. Southard},
  \bibinfo{title}{Perry's Chemical Engineers' Handbook}, \bibinfo{edition}{9th}
  ed., \bibinfo{publisher}{McGraw-Hill}, \bibinfo{address}{New York},
  \bibinfo{year}{2019}.

\end{thebibliography}


\end{document}
